\section{Постановка задачи}
\label{sec:task}

Целью лабораторной работы является изучение механизмов избирательного управления доступом: на уровне объектов файловой системы операционной системы Linux и на уровне ролей и привилегий в системе управления базами данных PostgreSQL.

В части~1 требуется разработать программу (набор сценариев командной оболочки), которая создаёт группы пользователей \texttt{group\_iit1}, \texttt{group\_iit2} и пользователей \texttt{iit11}, \texttt{iit12}, \texttt{iit21}, \texttt{iit22}, \texttt{iit3} с указанным членством в группах, а также предоставляет пользователю \texttt{iit21} административные привилегии; создаёт каталог \texttt{/home/pzs} и подкаталоги \texttt{pzs11}--\texttt{pzs15} с правами доступа только для владельца, только для группы, только для остальных, для всех пользователей и только для администратора (\texttt{root}); от имени пользователя \texttt{iit11} создаёт в доступных каталогах набор файлов \texttt{file11}--\texttt{file55} с комбинациями прав чтения, записи и исполнения; проверяет для указанных пользователей и суперпользователя возможность чтения, изменения и исполнения файлов, остановки процессов, запущенных по шаблону \texttt{filex5}, а также чтения содержимого каталогов, создания и удаления файлов в них; удаляет созданные файлы, каталоги, пользователей и группы.

В части~2 (вариант~\mdash  PostgreSQL) требуется установить СУБД, настроить механизм контроля доступа с ролями администратора (полный доступ), пользователя (CRUD) и гостя (только чтение), проверить корректность политики безопасности и удалить созданные объекты СУБД.
